% =====================================================================
% 5 CONCLUSÃO
% Elementos exigidos pela professora: conclusão (por objetivo), contribuições,
% limitações/pontos positivos e negativos, e trabalhos futuros de PESQUISA.
% Números da rodada R1 (cenário quente); conferir com a R2 quando a análise dela estiver pronta.
% \PENDENTE marca o que depende da R2, do custo de escrita corrigido e da iteração 2.
% =====================================================================

\section{CONCLUSÃO}\label{sec:conclusao}

Este trabalho teve como objetivo geral propor um \textit{framework} metodológico de quatro camadas para a otimização de consultas em SQL Server, tendo como critério de eficácia a redução mediana de pelo menos 20\% nas leituras lógicas de consultas com antipadrões de escrita. A hipótese foi confirmada para o conjunto das consultas: a redução mediana foi de 99,9\% nas 17 consultas principais (teste de Wilcoxon exato, $p < 0{,}001$), e 16 delas atingiram o limiar individualmente. O objetivo geral foi atingido, com a ressalva de que, na primeira volta, a prescrição não atingiu o limiar em uma consulta, S2, cuja redução foi de 6,4\%.

Quanto aos objetivos específicos, a revisão de oito trabalhos correlatos qualificados, quatro internacionais e quatro nacionais, na introdução, mostrou que as soluções atuam sobre a consulta já escrita e visam à automação, e que falta método que oriente o desenvolvedor durante a escrita. O catálogo reuniu quatro classes de antipadrões, com quatro ou cinco consultas cada, 19 no total, nas versões original e otimizada. O modelo do \textit{framework} definiu as quatro camadas, os critérios verificáveis de cada uma e o limiar de melhoria significativa, que o conjunto atingiu. O experimento antes e depois foi executado nas 19 consultas sobre o Stack Overflow 2010, com 20 repetições por consulta no cenário quente e 8 no frio. As reduções medianas das leituras lógicas, por classe, foram de 99,9\% em função aplicada à coluna filtrada, 63,1\% em subconsulta correlacionada, 100\% em \texttt{SELECT *} e 100\% em varredura por índice ausente. As reduções medianas do tempo de processamento e do tempo de execução, na segunda rodada, com a máquina sem outra carga, foram de 99,0\% nas duas métricas, e a classe de subconsulta correlacionada teve de novo o menor efeito, com 52,2\% e 51,6\%.

As contribuições do trabalho são o \textit{framework} de quatro camadas, o catálogo de antipadrões com consultas nas versões original e otimizada, o protocolo de medição, que reporta separadamente os cenários de cache frio e quente, inclui o custo de escrita e foi registrado antes da coleta, e a evidência experimental de que a reescrita sem índice não reduziu as leituras nas classes de função e de \texttt{SELECT *}, de que um índice que não cobre a consulta pode piorá-la em 2,06 vezes e de que a ordem das colunas do índice alterou as leituras em 75 vezes.

Como pontos positivos, a verificação por métricas objetivas mostrou que um índice pode piorar a consulta, e a regra de iteração do \textit{framework} identificou S2 como caso a revisar. Em S2 e S3, a prescrição manual indexou a tabela interna da junção, e o índice sugerido pela ferramenta nativa, no filtro externo, reduziu 89,6\% e 70,3\%, de modo que o método manual e a ferramenta automática se complementam. O custo de escrita dos índices do \textit{framework} foi de 1 a 3 vezes nos bytes de log, e o dos índices nativos para \texttt{SELECT *} foi de até 28 vezes. Na segunda volta do \textit{framework}, pós-hoc e não pré-registrada, S2 atingiu 95,9\% e S3, 99,2\%, ao se acrescentar o índice do filtro externo. Como limitações, destacam-se a emulação de arquitetura, que restringe a comparação absoluta de tempo de processamento e de tempo de execução, o uso de um único conjunto de dados e de um único sistema gerenciador, a ausência de concorrência entre usuários, o grau de paralelismo fixado em 1, a tabela maior que a memória do servidor e a aplicação do \textit{framework} pelo próprio autor.

Como trabalhos futuros de pesquisa, sugerem-se avaliar com desenvolvedores se a orientação do \textit{framework} altera a escrita das consultas, investigar a relação entre a seletividade das colunas e a magnitude do ganho, examinar a generalização do catálogo para outros conjuntos de dados, outros sistemas gerenciadores e cargas com usuários simultâneos, e estudar a combinação do método manual com a sugestão nativa de índices, que os resultados indicam ser complementares. \PENDENTE{validar com a orientadora os trabalhos futuros}
